\chapter{Anti-dictionnaire}

L'objectif de cette phase est de construire un anti-dictionnaire (une liste de mots non porteurs de sens à exclure du corpus avant indexation), en s'appuyant sur le coefficient TF-IDF pour les identifier automatiquement.

\section{Choix de l'unité documentaire}

On choisit \textbf{l'article} comme unité documentaire plutôt que le bulletin. Un bulletin regroupe des articles de rubriques très différentes, ce qui diluerait le signal TF-IDF en mélangeant des contextes hétérogènes. L'article, unité cohérente traitant d'un seul sujet, rend la fréquence d'un token bien plus informative. Ce choix est cohérent avec l'unité documentaire retenue pour la recherche.

\section{Calcul du TF-IDF}

Le calcul est réalisé par la classe \texttt{TFIDFProcessor}. Le TF est la fréquence brute du token dans le document, la fréquence normalisée n'étant pas nécessaire ici puisqu'on cherche à identifier les tokens les plus fréquents globalement. L'IDF est calculé selon la formule vue en cours :

\[
    \text{idf}_t = \log_{10}\!\left(\frac{N}{\text{df}_t}\right)
\]

Avant de calculer $\text{df}_t$, on déduplique les couples \texttt{(document\_id, token)} pour qu'un token apparaissant plusieurs fois dans un même document ne soit compté qu'une seule fois. On n'applique pas de lissage : un token présent dans tous les documents obtient un IDF de $\log_{10}(1) = 0$, ce qui annule directement son TF-IDF. Le TF-IDF final est le produit $\text{tf}_{t,d} \times \text{idf}_t$. Les tableaux contenant ces trois métriques pour chaque document du corpus sont sauvegardés dans des fichiers TSV séparés.
\section{Détermination de l'anti-dictionnaire}

La classe \texttt{AntiDict} identifie les stopwords en calculant le TF-IDF \textbf{moyen} de chaque token sur l'ensemble des documents où il apparaît. On retient la moyenne plutôt que la médiane ou le maximum : elle reflète le niveau global de discrimination des tokens et permet de prendre en compte les faiblesses récurrentes, contrairement au maximum qui dépend d’un seul pic, et à la médiane qui peut masquer certaines variations. 

Le seuil de $\text{tf-idf} \le 0.7$ à été retenue pour discriminer les stop-words des mots a garder. Cette valeur à été trouvée à partir de données, en les analysant avec le notebook \texttt{antidictionnaire.ipynb}.
 
Tout token dont la moyenne est inférieure ou égale au seuil est ajouté à l'anti-dictionnaire, stocké dans un \texttt{set} Python pour des recherches en $O(1)$ lors de l'execution.

La méthode \texttt{build\_sub\_table()} de \texttt{AntiDict} permet de l'avoir sous forme d'une table de substitution à deux colonnes (\texttt{token}, \texttt{sub}), où une chaîne vide \texttt{""} indique la suppression. Cette méthode est utile pour la classe \texttt{DataCleaner} qui se base sur ce format. Dans le 

\section{Génération du corpus sans stop word}

Le corpus sans stop word est déterminé avec le script \texttt{exercices-td/td2-determination\_antidict.py}, en particulier par la fonction \texttt{substitue()} qui applique la table de substitution de l'anti-dictionnaire au corpus. Ce script réponds à l'énoncé du TD, et permet d'avoir le fichier \texttt{corpus\_filtre.xml} nescessaire pour les prochains TD. Cependant, il n'est pas utilisé dans le programme finale, la logique ayant ensuite été déplacée dans la classe \texttt{DataCleaner} qui effectue toute la pipeline de nettoyage de données sans fichiers intermédiaires.  

